% ============================================================
% BAB 7 - ANTARMUKA EKSTERNAL
% ============================================================

\subsection{User Interface Requirements}

UI harus mendukung pekerjaan operasional yang berbeda per \textit{role}:
\textit{buyer} menelusuri dan membeli produk, \textit{seller} mengelola katalog
dan \textit{order} masuk, \textit{superadmin} memantau platform secara
keseluruhan. Karena itu halaman tidak hanya berfungsi sebagai \textit{form},
tetapi juga sebagai \textit{control surface} untuk status transaksi, stok, dan
kesehatan integrasi.

\begin{longtable}{|L{4cm}|L{2.5cm}|L{8cm}|}
    \caption{User Interface Requirements per Route}
    \label{tab:ui-req} \\
    \hline
    \textbf{UI / Route} & \textbf{Primary User} & \textbf{Requirement} \\
    \hline
    \endfirsthead
    \hline
    \textbf{UI / Route} & \textbf{Primary User} & \textbf{Requirement} \\
    \hline
    \endhead
    \hline
    \endfoot
    \texttt{/}
        & Semua / \textit{Guest}
        & \textit{Landing page} harus menampilkan katalog produk unggulan,
          \textit{banner} promo, dan akses \textit{login}/register. \\
    \hline
    \texttt{/auth/login} dan \texttt{/auth/register}
        & Semua
        & Halaman auth harus jelas membedakan opsi register sebagai \textit{buyer}
          atau \textit{seller}; tidak ada opsi \textit{superadmin}. \\
    \hline
    \texttt{/products}
        & \textit{Buyer} / \textit{Guest}
        & Halaman katalog harus mendukung filter kategori, \textit{search},
          \textit{sorting} harga, dan \textit{pagination}. \\
    \hline
    \texttt{/products/:id}
        & \textit{Buyer} / \textit{Guest}
        & Halaman detail produk harus menampilkan gambar, deskripsi, harga, stok,
          info toko, dan daftar ulasan. \\
    \hline
    \texttt{/cart}
        & \textit{Buyer}
        & Halaman keranjang harus menampilkan item, kuantitas (\textit{editable}),
          subtotal per item, dan tombol \textit{checkout}. \\
    \hline
    \texttt{/checkout}
        & \textit{Buyer}
        & Halaman \textit{checkout} harus menampilkan ringkasan \textit{order},
          kolom alamat pengiriman, \textit{input} \textit{voucher},
          \textit{fee breakdown} lengkap, dan tombol konfirmasi bayar. \\
    \hline
    \texttt{/orders} dan \texttt{/orders/:id}
        & \textit{Buyer}
        & Halaman \textit{order} harus menampilkan status terkini,
          \textit{tracking info}, dan opsi \textit{review} setelah
          \texttt{delivered}. \\
    \hline
    \texttt{/seller/dashboard}
        & \textit{Seller}
        & \textit{Dashboard} \textit{seller} harus menampilkan ringkasan produk
          aktif, stok rendah, dan \textit{order} masuk terbaru. \\
    \hline
    \texttt{/seller/products}
        & \textit{Seller}
        & Halaman manajemen produk harus mendukung CRUD produk dengan form
          validasi dan \textit{upload} gambar. \\
    \hline
    \texttt{/seller/orders}
        & \textit{Seller}
        & Halaman \textit{order} \textit{seller} menampilkan \textit{order} masuk
          dengan filter status dan detail item per \textit{order}. \\
    \hline
    \texttt{/admin}
        & \textit{Superadmin}
        & \textit{Dashboard} admin harus menampilkan metrik platform: total
          \textit{revenue}, \textit{order count}, \textit{user} aktif, dan grafik
          tren. \\
    \hline
    \texttt{/admin/users}
        & \textit{Superadmin}
        & Halaman manajemen \textit{user} menampilkan semua \textit{user} dengan
          aksi ban/aktifkan dan \textit{input} alasan. \\
    \hline
    \texttt{/admin/orders}
        & \textit{Superadmin}
        & Halaman \textit{order} admin menampilkan semua \textit{order} dengan
          kemampuan filter status dan \textit{update} manual. \\
    \hline
    \texttt{/admin/analytics}
        & \textit{Superadmin}
        & Halaman \textit{analytics} menampilkan produk terlaris,
          \textit{revenue} per \textit{seller}, distribusi kategori, dan tren
          transaksi. \\
    \hline
\end{longtable}

% ----
\subsection{API Requirements}

\begin{longtable}{|L{4.5cm}|L{1.5cm}|L{3.5cm}|L{5cm}|}
    \caption{Daftar API Endpoint PasarKita}
    \label{tab:api-req} \\
    \hline
    \textbf{Endpoint} & \textbf{Method} & \textbf{Auth} & \textbf{Deskripsi} \\
    \hline
    \endfirsthead
    \hline
    \textbf{Endpoint} & \textbf{Method} & \textbf{Auth} & \textbf{Deskripsi} \\
    \hline
    \endhead
    \hline
    \endfoot
    \texttt{/api/auth/register}
        & POST & ---
        & Registrasi \textit{user} baru; \textit{role} \texttt{buyer} atau
          \texttt{seller}. \\
    \hline
    \texttt{/api/auth/login}
        & POST & ---
        & \textit{Login}; mengembalikan JWT. \\
    \hline
    \texttt{/api/products}
        & GET & ---
        & \textit{Browse} semua produk aktif; mendukung filter dan
          \textit{pagination}. \\
    \hline
    \texttt{/api/products/:id}
        & GET & ---
        & Detail produk termasuk \textit{rating} agregat. \\
    \hline
    \texttt{/api/products}
        & POST & \textit{seller}
        & Tambah produk baru milik \textit{seller}. \\
    \hline
    \texttt{/api/products/:id}
        & PUT & \textit{seller} / \textit{superadmin}
        & Edit produk; \textit{seller} hanya produk milik sendiri. \\
    \hline
    \texttt{/api/products/:id}
        & DELETE & \textit{seller} / \textit{superadmin}
        & \textit{Soft delete} produk (\texttt{is\_active = false}). \\
    \hline
    \texttt{/api/checkout}
        & POST & \textit{buyer}
        & Proses \textit{checkout} lengkap termasuk \textit{payment request} ke
          SmartBank. \\
    \hline
    \texttt{/api/fee/calculate}
        & POST & ---
        & Simulasi kalkulasi \textit{fee} dan diskon tanpa transaksi. \\
    \hline
    \texttt{/api/orders}
        & GET & \textit{buyer} / \textit{seller} / \textit{superadmin}
        & Daftar \textit{order}; difilter berdasarkan \textit{role} yang
          \textit{request}. \\
    \hline
    \texttt{/api/orders/:id}
        & GET & \textit{buyer} / \textit{seller} / \textit{superadmin}
        & Detail \textit{order} termasuk \textit{items} dan status
          \textit{history}. \\
    \hline
    \texttt{/api/orders/:id/status}
        & PATCH & \textit{superadmin}
        & \textit{Update} status \textit{order} manual. \\
    \hline
    \texttt{/api/orders/:id/ratings}
        & POST & \textit{buyer}
        & \textit{Submit} \textit{rating} setelah \textit{order} \texttt{delivered}. \\
    \hline
    \texttt{/api/admin/users}
        & GET & \textit{superadmin}
        & Semua \textit{user} dengan filter dan \textit{pagination}. \\
    \hline
    \texttt{/api/admin/users/:id/status}
        & PATCH & \textit{superadmin}
        & Ban atau aktifkan \textit{user}. \\
    \hline
    \texttt{/api/admin/analytics}
        & GET & \textit{superadmin}
        & Metrik platform: \textit{revenue}, \textit{order}, produk terlaris. \\
    \hline
    \texttt{/api/seller/profile}
        & GET / PUT & \textit{seller}
        & Baca dan \textit{update} profil toko \textit{seller}. \\
    \hline
    \texttt{/api/notifications}
        & GET & semua \textit{role}
        & Daftar notifikasi milik \textit{user} yang sedang \textit{login}. \\
    \hline
    \texttt{/api/notifications/:id/read}
        & PATCH & semua \textit{role}
        & Tandai notifikasi sudah dibaca. \\
    \hline
\end{longtable}

% ----
\subsection{Data Exchange Contract}

Kontrak JSON pada \textit{checkout request} dan \textit{response} harus
diperlakukan sebagai \textit{integration contract} dengan API Gateway dan
\textit{service partner}. Perubahan nama \textit{field} atau tipe data dapat
memutus integrasi dengan SmartBank atau LogistiKita.

\begin{lstlisting}[language={},caption={POST /api/checkout: Request Body},
    label={lst:checkout-req}]
// POST /api/checkout -- Request Body
{
  "items": [
    { "product_id": "uuid", "qty": 2 }
  ],
  "shipping_address": "Jl. Contoh No. 1, Bandung",
  "voucher_code": "UMKM2026",
  "idempotency_key": "uuid-v4"
}
\end{lstlisting}

\begin{lstlisting}[language={},caption={POST /api/checkout: Response (sukses)},
    label={lst:checkout-res}]
// POST /api/checkout -- Response (sukses)
{
  "order_id": "uuid",
  "status": "paid",
  "subtotal": 100000,
  "fee_marketplace": 2000,
  "voucher_discount": 5000,
  "total": 97000,
  "transaction_id": "SB-TXN-XXXX",
  "tracking_id": "LK-TRACK-XXXX"
}
\end{lstlisting}
